\chapter[\emj{🐢}~Fast \& Slow Pointers (Tortuga y Liebre)]{\emj{🐢}~Fast \& Slow Pointers (Tortuga y Liebre)}

\begin{dsaquees}

Dos corredores en una pista circular 🏃: uno lento (tortuga, avanza 1
paso) y uno rápido (liebre, avanza 2 pasos). Si la pista tiene un bucle,
tarde o temprano el rápido le da la vuelta y ALCANZA al lento por
detrás. Si la pista termina (hay una meta), el rápido llega al final y
nunca se encuentran.

Ese simple truco resuelve un montón de problemas de listas enlazadas:
detectar ciclos, encontrar el nodo del medio y hallar dónde empieza el
bucle, todo sin usar memoria extra.

\end{dsaquees}

\begin{dsaotraforma}

La técnica Fast \& Slow (algoritmo de la tortuga y la liebre de Floyd)
usa dos punteros que recorren la lista a velocidades distintas: uno
avanza de 1 en 1 y otro de 2 en 2.

Sus tres usos clásicos:
\begin{itemize}
\item \textbf{Detectar ciclo}: si en algún momento \verb|fast == slow|,
      hay ciclo. Si \verb|fast| llega a NULL, no lo hay.
\item \textbf{Encontrar el medio}: cuando \verb|fast| llega al final,
      \verb|slow| está justo a la mitad.
\item \textbf{Inicio del ciclo}: tras el encuentro, mueve un puntero al
      \verb|head| y avanza ambos de 1 en 1; se cruzan en el inicio del
      bucle.
\end{itemize}

Todo en $O(n)$ tiempo y $O(1)$ espacio: ahí está su valor.

\end{dsaotraforma}

\begin{dsadiagrama}
\begin{verbatim}
Lista con ciclo:

  1 -> 2 -> 3 -> 4 -> 5
            ^         │
            └─────────┘

slow (1 paso):  1  2  3  4  5  3  4 ...
fast (2 pasos): 1  3  5  4  3  5  4 ...
                      encuentro en el nodo 4  => HAY CICLO

Encontrar el medio (lista sin ciclo):
  1 -> 2 -> 3 -> 4 -> 5 -> NULL
fast llega a NULL  =>  slow queda en 3 (el medio)
\end{verbatim}
\end{dsadiagrama}

\begin{lstlisting}
DETECTAR CICLO
- slow = head, fast = head.
- Avanza slow 1 y fast 2 mientras fast y fast->next existan.
- Si slow == fast en algún punto -> hay ciclo.

INICIO DEL CICLO (después del encuentro)
- Deja un puntero en el punto de encuentro, mueve otro al head.
- Avanza ambos de 1 en 1; donde se crucen es el inicio del ciclo.

📝 PSEUDOCÓDIGO (detectar + inicio de ciclo)
──────────────────────────────────────────────────────────────

función detectarCiclo(head):
    slow = head; fast = head
    MIENTRAS fast Y fast.next:
        slow = slow.next
        fast = fast.next.next
        SI slow == fast:                 // hay ciclo
            p = head
            MIENTRAS p != slow:
                p = p.next; slow = slow.next
            return p                      // inicio del ciclo
    return NULL                           // no hay ciclo

💻 CÓDIGO C++
──────────────────────────────────────────────────────────────

struct Node { int val; Node* next; };

// ¿Tiene ciclo?
bool hasCycle(Node* head) {
    Node *slow = head, *fast = head;
    while (fast && fast->next) {
        slow = slow->next;
        fast = fast->next->next;
        if (slow == fast) return true;   // se alcanzaron
    }
    return false;                        // fast llegó a NULL
}

// Nodo del medio
Node* middleNode(Node* head) {
    Node *slow = head, *fast = head;
    while (fast && fast->next) {
        slow = slow->next;
        fast = fast->next->next;
    }
    return slow;                         // mitad exacta (o 2da mitad si par)
}

⏱️ COMPLEJIDAD (Big-O)
──────────────────────────────────────────────────────────────

  Problema          │ Tiempo  │ Espacio
  ──────────────────┼─────────┼──────────
  Detectar ciclo    │ O(n)    │ O(1)
  Encontrar medio   │ O(n)    │ O(1)
  Inicio del ciclo  │ O(n)    │ O(1)

* La gran ventaja es el O(1) de memoria (sin HashSet de nodos vistos).
\end{lstlisting}

\begin{lstlisting}
💡 Otros usos del patrón:
   - Palíndromo de lista: encuentra el medio, invierte la 2da mitad y
     compara.
   - Happy Number: aplica fast/slow sobre la secuencia de sumas de
     cuadrados de dígitos para detectar el ciclo.
   - Eliminar el N-ésimo desde el final: separa dos punteros por N
     nodos y avánzalos juntos.
\end{lstlisting}

\begin{dsaentrevistas}

✅ "Linked List Cycle" (LeetCode 141/142): 141 pide detectar, 142 pide
   el nodo donde empieza. Ambos son Floyd puro.

💡 ¿Por qué funciona lo del inicio? La distancia del head al inicio del
   ciclo es igual a la distancia del punto de encuentro al inicio del
   ciclo (avanzando en el bucle). Sabértelo justifica el algoritmo.

⚠️ Cuida los NULL: siempre valida \verb|fast != NULL| Y
   \verb|fast->next != NULL| antes de hacer \verb|fast->next->next|, o
   provocas un segfault.

🔑 Alternativa con HashSet: guardar nodos visitados también detecta
   ciclos, pero usa O(n) memoria. Floyd es superior por su O(1).

\end{dsaentrevistas}

\begin{dsaresumen}

• Dos punteros: lento (1 paso) y rápido (2 pasos).
• Se encuentran = hay ciclo; fast llega a NULL = no hay.
• fast al final -> slow queda en el medio.
• Inicio del ciclo: un puntero al head, avanzar ambos de 1 en 1.
• O(n) tiempo, O(1) espacio (mejor que HashSet).
• Clásicos: cycle detection, medio, palíndromo, happy number.
════════════════════════════════════════════════════════════════

\end{dsaresumen}
